\subsection{An inverse quintic construction and actual bad transferred images}

\paragraph{Finding A40 (inverse dihedral construction in characteristic two).}
Let $E$ be an ordinary smooth projective genus-two curve over an
algebraically closed field $k$ of characteristic two. Suppose its
canonically polarized Jacobian $J$ has an integral Rosati-symmetric
endomorphism
\[
                    T^2+T-1=0,\qquad T^\dagger=T.
 \tag{ID.1}
\]
There is a smooth connected genus-six curve $D$ with an action of
the dihedral group of order ten such that
\[
 \begin{gathered}
 D/\langle\sigma\rangle=C,\quad g(C)=2,\quad
 D\longrightarrow C\ \hbox{is etale of degree five},\\
 D/\langle\tau_1\rangle\simeq E\simeq D/\langle\tau_2\rangle.
 \end{gathered}
 \tag{ID.2}
\]
where $\sigma=\tau_1\tau_2$ has order five. Each $\tau_i$
descends to the hyperelliptic involution of $C$.

If the initial data are over a finite field, the construction and
all these maps are defined after a finite constant extension.
Every prescribed geometric point $R$ of $E$ is then the image of
a rational lift of a point $P$ on $C$ for a constant, $P$-split
cyclic degree-five cover, after one further finite extension.
No assertion that this $C$ has two-rank one, or that this $P$ is
Weierstrass, is part of the statement.

\paragraph{Correspondence background.}
Kumar--Mukamel describe the passage from a Jacobian endomorphism
to an effective correspondence and give a genus-six $(2,2)$
real-multiplication example over a number field
\cite[Sections~2--3, pp.~5--7]{KumarMukamel2016}.
The proof here establishes integrality and generic smoothness
for each fixed ordinary characteristic-two curve, and then
proves the unramified dihedral and pointed finite-field
assertions. No priority claim for the general correspondence
method is made.

\begin{proof}
\emph{1. The residual theta correspondence and its exact numerical class.}
Choose a degree-one theta characteristic $B$ on $E$, so $2B\sim K_E$,
and write
\[
             z_P=[P-B],\qquad \Theta=\{z_P:P\in E\}\subset J.
 \tag{ID.3}
\]
This is a smooth symmetric theta divisor: $z_{\iota P}=-z_P$,
where $\iota$ is the hyperelliptic involution. The construction
does not require this choice to be the Raynaud characteristic.
Equation (ID.1) makes $T$ an automorphism, with inverse $T+1$.

We first record the intersection numbers, including their
normalization. The characteristic polynomial of $T$ on a
four-dimensional rational Tate module is $(X^2+X-1)^2$.
Indeed it has rational coefficients and its only possible
eigenvalues are the two conjugate roots of the irreducible
quadratic in (ID.1), with equal multiplicities. Thus
$\operatorname{Tr}(T)=-2$, and
\[
 \operatorname{Tr}(T^\dagger T)=6,
 \qquad
 \operatorname{Tr}((T^{-1})^\dagger T^{-1})=6,
 \qquad \Theta\cdot T\Theta=3.
 \tag{ID.4}
\]
Here $\Theta^2=2$ and the last equality is the Rosati trace
formula $\Theta\cdot a^*\Theta=\operatorname{Tr}(a^\dagger a)/2$,
applied to $a=T^{-1}$; see Milne, Theorem~17.3 and its proof
\cite[pp.~35--36]{MilneAV}. In particular $T\Theta$ cannot
be a translate of $\Theta$, since their self-intersections
are two whereas their mutual intersection is three.

For $t\in E$ define
\[
 F_t:E\times E\longrightarrow J,
       \qquad F_t(P,Q)=T(z_P-z_t)-z_Q,
 \qquad V_t=\{t\}\times E.
 \tag{ID.5}
\]
The map $F_t$ is finite of degree three. Its fibres are the
intersections of $T\Theta$ with translates of $\Theta$, so
they are finite by (ID.4); their generic intersection degree
is three. It is also generically separable: the tangent
directions $dT(T_P E)$ and $T_Q E$ are generically distinct,
because the canonical Gauss map of a genus-two Abel curve
is nonconstant. Being finite between smooth surfaces, $F_t$
is flat by the regular/Cohen--Macaulay flatness criterion
\cite[Lemma~10.128.1, Tag~\texttt{00R4}]{Stacks00R4}.
Put
\[
                  D_t=F_t^*\Theta-V_t.
 \tag{ID.6}
\]
This is an effective divisor. The component $V_t$ occurs in
$F_t^*\Theta$ with multiplicity one: along its generic point
the normal derivative is the pairing of the nonzero vector
$dT(T_tE)$ with the conormal of $\Theta$ at $-z_Q$, and is
not identically zero as $Q$ varies.

There is no other vertical component. Such a component over
$P$ would give $\Theta+T(z_P-z_t)=\Theta$, whose translation
is zero because a principal theta divisor has trivial
translation stabilizer. Then $P=t$. There is no horizontal
component either, since that would make $T\Theta$ a translate
of $\Theta$. The two projection degrees of $D_t$ are both two:
the original degrees of $F_t^*\Theta$ are two over the first
factor and three over the second, and subtracting $V_t$
decreases only the second degree.

The same construction determines the sign of its Jacobian
correspondence. For $P\ne t$, its two points in the second
factor, counted with their multiplicities, form the unique
effective divisor
\[
                   Q+R\sim K_E+T([P-t]).
 \tag{ID.7}
\]
Indeed $F_t(P,Q)=z_R$ is precisely this equality, and the
degree-two line bundle has a unique effective divisor when
$T([P-t])\ne0$, by Riemann--Roch. This also identifies the
intersection multiplicities: away from the canonical class
the Abel map $E^{(2)}\to\operatorname{Pic}^2(E)$ is an
isomorphism. Therefore the induced homomorphism from the
first Jacobian to the second is $+T$, not $-T$.

Finally, finite flat pullback of intersection numbers gives
$(F_t^*\Theta)^2=3\Theta^2=6$; also
$F_t^*\Theta\cdot V_t=2$ and $V_t^2=0$.
Consequently
\[
 D_t^2=2,\qquad
 D_t\cdot K_{E\times E}=8,\qquad p_a(D_t)=6.
 \tag{ID.8}
\]
The pullback formula is Milne, Lemma~8.3
\cite[p.~13]{MilneAV}; the last equality is adjunction on
the smooth product surface.

\emph{2. The residual divisor is integral.}
Every irreducible component has positive degree over both
factors. A divisor of bidegree $(2,2)$ which is reducible or
nonreduced must therefore be the sum of two graphs of
automorphisms $\varphi,\psi$ of $E$, allowing the two graphs
to coincide. A component of bidegree $(1,1)$ is a graph
because its normalization maps with degree one to each
smooth curve $E$.

Let $u,v$ be the induced automorphisms of $J$. They preserve
the canonical polarization, so $u^\dagger=u^{-1}$ and
$v^\dagger=v^{-1}$. By (ID.7), $T=u+v$. With
$w=u^{-1}v$ it follows that
\[
              T^2=T^\dagger T=2+w+w^{-1},
       \qquad w+w^{-1}=-1-T.
 \tag{ID.9}
\]
The latter element satisfies $S^2+S-1=0$. Multiplication
by $w^2$ therefore gives
\[
                  1+w+w^2+w^3+w^4=0.
 \tag{ID.10}
\]
Hence $w$ has exact order five; it cannot be the identity
since multiplication by five is not zero on $J$.
Automorphisms of a genus-at-least-two curve act faithfully
on its Jacobian. This follows, for example, from the
uniqueness argument in Milne's Theorem~12.1(b)
\cite[pp.~38--39]{MilneJV}.
Thus $\varphi^{-1}\psi$ would be an automorphism of $E$
of order five.

This is impossible for an ordinary genus-two curve in
characteristic two. Its separable hyperelliptic map has
exactly three geometric branch points. Every automorphism
descends to the projective line and permutes these three
points; the kernel is the hyperelliptic group of order two,
and the induced group embeds in $S_3$. Its order divides
twelve. This contradiction proves that $D_t$ is integral
and reduced, for every $t$.

\emph{3. All but finitely many parameters give a smooth curve.}
Here the characteristic-two argument is essential. Let
\[
                \gamma:E\longrightarrow\mathbf P(T_0J)
 \tag{ID.11}
\]
send $P$ to the tangent line of its Abel image translated
to the origin. Pullback of invariant differentials identifies
this map with the canonical map of $E$. Thus $\gamma$ is the
separable hyperelliptic double map, and its geometric fibres
are $\{P,\iota P\}$. Its ramification points are precisely
the three Weierstrass points. Let $\overline T$ be the
projective automorphism induced by the invertible differential
$dT$.

Consider the reduced critical curve
\[
 Z=\{(P,Q):\gamma(Q)=\overline T\gamma(P)\}
             \subset E\times E,
 \qquad h(P,Q)=Tz_P-2z_Q.
 \tag{ID.12}
\]
On the normalization of every irreducible component $Z_0$,
both projections to $E$ are nonconstant and generically
separable. Indeed this is a fibre product of two finite
separable double maps, and the assertion holds off their
finite branch sets. Since $d[2]=0$ in characteristic two,
\[
                         dh=dT\,dz_P\ne0
                         \quad\hbox{generically on }Z_0.
 \tag{ID.13}
\]
In particular $h$ is nonconstant on each such component.

No image $h(Z_0)$ is contained in $T\Theta$. Otherwise
there would be a nonconstant map $s:Z_0\to E$ with
$Tz_s=h(P,Q)$. Differentiating and using (ID.13) gives
$\gamma(s)=\gamma(P)$ generically. The two points in a
canonical fibre imply $s=P$ or $s=\iota P$ on this
irreducible component. The first possibility would make
$2z_Q=0$ identically, contradicting the nonconstant second
projection. The second gives
$2(Tz_P-z_Q)=0$. The map $Tz_P-z_Q$ would then be constant,
because its image lies in the finite group scheme $J[2]$
and the source is reduced and connected. Surjectivity of
both projections would make $T\Theta$ a translate of
$\Theta$, contrary to (ID.4). Thus each curve $h(Z_0)$
has only finitely many intersections with $T\Theta$.

Suppose $(P,Q)$ with $P\ne t$ is singular on $F_t^*\Theta$,
and write $F_t(P,Q)=z_R$. The differential criterion for
this pullback of a smooth divisor is
\[
                    \gamma(R)=\gamma(Q)=\overline T\gamma(P).
 \tag{ID.14}
\]
Thus $R=Q$ or $R=\iota Q$. The latter would give
$T(z_P-z_t)=0$, hence $P=t$, which was excluded. In the
former case $(P,Q)\in Z$ and $h(P,Q)=Tz_t$.
The previous paragraph shows that this can happen for only
finitely many parameters $t$. Outside that finite set,
$D_t$ is smooth away from $V_t$.

It remains to check the intersections with $V_t$, rather
than simply discard this fixed component. Near $(t,Q)$
take a local parameter $a$ on the first factor. A local
equation of $F_t^*\Theta$ has the form $aG(a,Q)$, and $G$
defines the residual divisor. Up to a local unit,
$G(0,Q)$ is the pairing of the conormal of $\Theta$ at
$-z_Q$ with $dT(T_tE)$. Its zeros are exactly
\[
                         \gamma(Q)=\overline T\gamma(t).
 \tag{ID.15}
\]
At a non-Weierstrass $Q$ this zero is simple, since
$\gamma$ is etale there. Hence the derivative of $G$ in
the second-factor parameter is nonzero and the residual
curve is smooth at this point. Only finitely many $t$
have $\overline T\gamma(t)$ equal to a branch value of
$\gamma$ (at most six geometric points). Excluding them
as well proves that $D=D_t$ is smooth for every $t$ in a
nonempty open subset of $E$. By (ID.8), it has genus six.

\emph{4. The two involutions give the unramified dihedral cover.}
Let $v_1,v_2:D\to E$ be the projections. Both have degree
two and are separable. A purely inseparable degree-two map
between smooth curves over this perfect field would be a
relative Frobenius followed by an isomorphism and would
preserve the genus, whereas here the genera are six and two.
Let $\tau_i$ be the involution of $v_i$, and put
\[
 U_i=v_i^*:J\longrightarrow J_D,
       \qquad N_i=v_{i,*}:J_D\longrightarrow J.
 \tag{ID.16}
\]
Pullback and norm are adjoint for the canonical principal
polarizations. Thus (ID.7), Rosati symmetry, and the degree
two identities give
\[
 N_iU_i=2,\quad U_iN_i=1+\tau_i,\quad
 N_2U_1=T=N_1U_2.
 \tag{ID.17}
\]
The map $U=(U_1,U_2):J^2\to J_D$ has finite kernel:
its adjoint composite is
\[
 NU=\begin{pmatrix}2&T\\T&2\end{pmatrix},
 \qquad
 (4-T^2)=(3+T),\qquad (3+T)(2-T)=5.
 \tag{ID.18}
\]
The adjugate shows that this matrix is an isogeny, since
multiplication by five is an isogeny. Its image $A$ has
dimension four. Its polarized complementary abelian
subvariety $B_D=\ker(N_1,N_2)^0$ has dimension two, and
$J_D$ is isogenous to $A\times B_D$.

Under the isogeny $J^2\to A$, the two involutions and their
product are represented by
\[
 \tau_1=\begin{pmatrix}1&T\\0&-1\end{pmatrix},\quad
 \tau_2=\begin{pmatrix}-1&0\\T&1\end{pmatrix},\quad
 \sigma=\tau_1\tau_2=
            \begin{pmatrix}-T&T\\-T&-1\end{pmatrix}.
 \tag{ID.19}
\]
The last matrix has determinant one and trace $-1-T$.
Therefore $\sigma+\sigma^{-1}=-1-T$ and, exactly as in
(ID.10), $\Phi_5(\sigma)=0$ on $A$.
Also $\sigma^2$ has matrix
$\left(\begin{smallmatrix}0&-1\\1&T\end{smallmatrix}\right)$,
so $N_1\sigma^2U_1=T$: choosing the squared deck generator
recovers the originally specified real-multiplication
endomorphism, rather than its conjugate $-1-T$.
On $B_D$, both $\tau_i$ act as $-1$ by (ID.17), so
$\sigma$ acts as $1$. It follows that $\sigma^5=1$ on
$J_D$ and that its action is not the identity. Faithfulness
of the action of curve automorphisms on the Jacobian gives
an automorphism of $D$ of exact order five. Together with
$\tau_1$ it generates the dihedral group, since
$\tau_1\sigma\tau_1=\sigma^{-1}$.

The $\sigma$-invariant part of the rational Tate module of
$J_D$ has dimension four: it is the part belonging to
$B_D$, with none from $A$ because $\Phi_5(1)=5\ne0$.
For the quotient $C=D/\langle\sigma\rangle$, pullback
identifies its Jacobian up to isogeny with this invariant
abelian subvariety; thus $g(C)=2$.
Riemann--Hurwitz now gives
\[
                 2g(D)-2=10=5(2g(C)-2),
 \tag{ID.20}
\]
so the different is zero and the degree-five map is etale.
The quotients by the two reflections are exactly the two
projection targets $E$. On $J_C\sim B_D$, either reflection
acts as $-1$. Its quotient has genus zero, so it is the
hyperelliptic involution of $C$. This proves (ID.2).

\emph{5. Finite-field descent and an arbitrary prescribed image.}
Suppose $E$ and the original data are over a finite field.
The theta characteristic, a point $t$ in the nonempty
smoothness open, the residual divisor and all its maps have
finitely many algebraic coefficients. They are defined over
one finite constant extension $k_1$; the two degree-two
deck involutions and their product are then defined over
that same field. Given any $R\in E(\overline{k_1})$, choose
$P'\in v_1^{-1}(R)$ and a finite extension $k_2/k_1$ over
which $R$ and $P'$ are rational. Put $P=\pi(P')\in C(k_2)$.
The entire fibre of the etale map $\pi:D\to C$ above $P$
is $\{\sigma^jP':0\le j<5\}$, so it is split over $k_2$.
Its deck group is constant. This proves the pointed
assertion without identifying $R$ with the auxiliary
parameter $t$ and without prescribing the rank of $C$.
\end{proof}

\paragraph{Consequences for the canonical forcing attempt.}
Apply A40 to the ordinary integral-$\mathcal O_5$ curve of
A13 and A38. Choose one of its $126$ non-Weierstrass
$\mathbf F_{128}$-rational points, for example (BC.7).
After the finite extensions just constructed, this is an
actual image of a split lift in a quintic transfer with an
ordinary reflection quotient. Its fixed canonical class
still has exact order $4159$, and
$\operatorname{ord}_{4159}(2)=2079$ is odd. Hence even among
actual prescribed images of quintic transfers from genus-two
curves, the A29 canonical-prime hypothesis need not hold.

This does not assert that both reflection images associated
to one lift fail, that every possible quintic cover fails,
or that the rank-one-source/Weierstrass-point subcase has
been refuted. It closes the broad image-membership gap left
in A38. The geometric order-$4159$ theta witness of A13
still provides point separation; no section collision is
constructed here. The proof is unconditional and does not
use a prime-density assertion.

The next finding A41 supplies the stronger simultaneous failure at all
five reflection images for a fixed-field member of this inverse family.
It still leaves a suitable choice of a different cover open.
